% Pista geo-25s-q4b · Geometria
% Procedència: PAU Catalunya 2025 · Convocatòria de setembre · Sèrie 3 · Exercici 4 — Opció B
\documentclass[11pt,a4paper]{article}
\usepackage{pista}
\pistainfo{Catalunya 2025}{Convocatòria de setembre}{Sèrie 3}

\begin{document}

\section*{\textcolor{blauPista}{Exercici 4 --- Opció B}}

\noindent Considereu el pla $\pi: 2x - y + z = 5$ i el punt $P = (0, 1, 3)$.

\subsection*{\textit{a)} \normalfont Comproveu que la distància del punt $P$ al pla $\pi$ és $\dfrac{\sqrt{6}}{2}$. \hfill\textnormal{[0,5~punts]}}

\pista{Aplica directament la fórmula de la distància d'un punt a un pla. Si el resultat et surt amb una arrel al denominador, racionalitza'l per veure que coincideix amb $\dfrac{\sqrt{6}}{2}$.}

\subsection*{\textit{b)} \normalfont Trobeu l'equació general d'un pla $\pi_{1}$ paral·lel a $\pi$ i que passi pel punt $P$. Quina és la distància entre $\pi_{1}$ i $\pi$? \hfill\textnormal{[0,75~punts]}}

\pista{Dos plans paral·lels tenen el mateix vector normal. Per tant, l'equació de $\pi_{1}$ és $2x - y + z + D = 0$, on $D$ és un terme independent que hauràs de determinar. Imposa que el pla passi per $P$ i així podràs trobar el valor de $D$. Per poder trobar la distància entre els plans $\pi$ i $\pi_{1}$, observa que ja l'has calculada a l'apartat a); justifica per què: quin punt de $\pi_{1}$ coneixes?}

\subsection*{\textit{c)} \normalfont Trobeu l'equació general d'un segon pla $\pi_{2}$, diferent de $\pi_{1}$, que és paral·lel a $\pi$ i que estigui a una distància $\dfrac{\sqrt{6}}{2}$ de $\pi$. \hfill\textnormal{[1,25~punts]}}

\pista{Com a l'apartat anterior, $\pi_{2}$ té la forma $2x - y + z + D = 0$. Imposa que la distància entre $\pi$ i $\pi_{2}$ sigui $\dfrac{\sqrt{6}}{2}$ (per exemple, pren un punt qualsevol de $\pi$ i calcula'n la distància a $\pi_{2}$): obtindràs una equació amb un valor absolut, que tindrà dues solucions. Una d'elles correspon a $\pi_{1}$; l'altra és, precisament, la que et convé.}

\end{document}
